\chapter{PHOTOGRAPHIC AND DOCUMENTARY EVIDENCE REGISTER}
\PUAddAppendixEntry{PHOTOGRAPHIC AND DOCUMENTARY EVIDENCE REGISTER}
\label{app:evidenceregister}
Photographs are supplementary to physical and documentary evidence. No unverified or automatically selected photograph is presented here as a completed rework case. A photograph should be identified by project, element, construction activity, date and permission to publish. Descriptions should remain neutral; the mere appearance of a crack or damp patch does not establish causation or corrective work.
\begin{longtable}{@{}L{0.14\textwidth}L{0.31\textwidth}L{0.43\textwidth}@{}}
\caption{Three current photograph and documentary categories}\label{tab:photoregister}\\
\toprule Code & Relevant subjects & Evidence required for case verification\\ \midrule\endfirsthead
\toprule Code & Relevant subjects & Evidence required for case verification\\ \midrule\endhead
Waterproofing & Bathroom floors, wall returns, drains, thresholds, terrace slopes, scaffolding penetrations, membrane substrates. & Approved waterproofing detail, original installation, source testing or inspection, corrective repair and restored finish documentation.\\
Concrete pavement & Building-associated PCC pavement panels, joints, edges, subbase, surface cracks and curing arrangements. & Scope confirmation, design, construction dates, slab/joint/support records, crack mapping, repair or recasting instruction and completion record.\\
Masonry/plaster & Perpendicular masonry walls, butt/bonded/connected junctions, openings, bands, brick moisture condition, curing and plaster cracks. & Approved masonry details, substrate condition, dated observations, inspection/repair instruction and any actual plaster or wall correction.\\
\bottomrule
\end{longtable}

\section{Photograph Admission Procedure}
An image is associated with the applicable category only after confirming its source, construction context and consent. The reviewer then separately checks whether original completed work and the corrective intervention are documented. Without both, the record remains a construction-quality illustration or an unverified observation; it must not be counted as rework. Historic photographs from previous candidate-screening exercises do not gain scientific status merely by being stored in the project folder.

\section{Records Still Required}
The manuscript now includes 25 photographs showing construction and repair-related conditions; their links to dated repair records and cost measurements are still to be assessed individually. Project permissions, dated source photographs, field notes, original-language practitioner documentation, photograph-to-site attribution, corrective-action and programme/cost records are pending. Missing information should be obtained from the actual relevant projects instead of being inferred from illustrations or literature.
